\section{Local collision constraints and a window of actual returns}
\label{sec:v37-return-window}
The occupation variable in the preceding section is exactly the missing
integer in an actual-return description at a prescribed collision time.
We now use this fact as an identity of measures. There is no random-time
approximation in the identity and no replacement of a rare event.

For $m\ge1$ define a finite measure on $\R$ by
\begin{equation}\label{eq:v37-return-window-measure}
 \mathfrak R_{m,R}^{k,t,J}
 =m^{3/2}\sum_{n=1}^{m}
 \nu_R^*\{K_{n,R}=k,\ N_{n,R}=m,\ T_{n,R}-t\in J\}
                         \delta_{(n-mc)/\sqrt m}.
\end{equation}
Here $J$ is a fixed bounded roof interval of positive length, and
$c=c_*$. Every term concerns the same four-coordinate physical return
record as Theorem~\ref{thm:LLT}. In particular $n$ is a genuine number
of returns, not a deterministic substitute for a collision count.

\begin{proposition}[Exact return and occupation disintegration]
\label{prop:v37-exact-return-disintegration}
For every bounded measurable $F:\R\to\mathbb C$,
\begin{align}\label{eq:v37-exact-return-disintegration}
 \int F\,d\mathfrak R_{m,R}^{k,t,J}
 =\frac{m^{3/2}}c\int_M\eta_R(x)\eta_R(T_R^mx)
 \one_{\{K^{\rm c}_{m,R}=k,\ S_{m,R}-t\in J\}}
 F\left(\frac{A_{m,R}-mc}{\sqrt m}\right)d\nu(x).
\end{align}
\end{proposition}
\begin{proof}
Discard the invariant null set of singular trajectories and undefined
returns. Start at $x\in Y_R^*$. If $T_R^mx\in Y_R^*$, list the section
visits at collision times $0=j_0<j_1<\cdots<j_n=m$.
There are exactly $n$ visits in $[0,m)$, so $A_{m,R}(x)=n$ and
$N_{n,R}(x)=m$. Conversely $N_{n,R}(x)=m$ implies this list and this
occupation count. Thus
\[
 \{x\in Y_R^*,\ T_R^mx\in Y_R^*,\ A_{m,R}=n\}
       =\{x\in Y_R^*,\ N_{n,R}=m\}.
\]
On that set additivity gives $K_{n,R}=K^{\rm c}_{m,R}$ and
$T_{n,R}=S_{m,R}$. The sets are disjoint as $n$ varies and exhaust
the terminal-section event. Since at most one visit occurs at each
collision, $1\le n\le m$. Integrate and use
$d\nu_R^*=c^{-1}\eta_R\,d\nu$. This proves the identity, including
both endpoint conventions.
\end{proof}

\begin{theorem}[A local limit in a diffusive actual-return window]
\label{thm:v37-return-window}
Let $\Sigma_R,\beta_R,\sigma_{A,R}^2$ be as in
\eqref{eq:v37-schur-data}, and set
$z=z_{m,R}(k,t)=(k/\sqrt m,(t-m\bar\tau_R)/\sqrt m)$.
For every fixed bounded continuous $F$, uniformly in $R\in I$ and
$|z|\le M$,
\begin{equation}\label{eq:v37-return-window-limit}
 \int F\,d\mathfrak R_{m,R}^{k,t,J}
  -c|J|g_{\Sigma_R}(z)
       \int_\R F(y)g_{\sigma_{A,R}^2}(y-\beta_Rz)\,dy
                       \longrightarrow0.
\end{equation}
The formula holds also when $F$ is the indicator of a fixed interval,
including a half-line. For each fixed $a<b$, the unscaled probability
sum in \eqref{eq:v37-return-window-measure} over
$a\le(n-mc)/\sqrt m\le b$ is at least $c_{M,a,b,J}m^{-3/2}$ for
all sufficiently large $m$, with $c_{M,a,b,J}>0$ uniform in $R$.
\end{theorem}
\begin{proof}
Apply Theorem~\ref{thm:v37-occupation-local-central} with
$a_R=d_R^{\rm end}=\eta_R$. These are admissible multipliers by
Lemma~\ref{lem:v37-section-multiplier}. Their product of means is
$c^2$, while the exact return disintegration has the factor $c^{-1}$.
This gives precisely the amplitude $c$ in
\eqref{eq:v37-return-window-limit}.
The normal law has no atoms, so fixed interval endpoints are continuity
sets. Continuity, uniform ellipticity and compactness give a positive
lower bound for
\[
 c|J|g_{\Sigma_R}(z)
       \int_a^b g_{\sigma_{A,R}^2}(y-\beta_Rz)\,dy
\]
on $R\in I$, $|z|\le M$. The uniform error tends to zero, proving the
last assertion. The measure is positive throughout the argument.
\end{proof}

\subsection{Agreement with the four-coordinate Gaussian normalization}
The new amplitude is compatible with the covariance of the original
induced record, rather than a newly chosen covariance.
From \eqref{eq:v37-covariance-change},
\begin{equation}\label{eq:v37-return-covariance-jacobian}
 \det\Omega_R=c^6\det D_R,\qquad
 L_R^{-1}(z_1,z_2,z_3,y)
       =(z_1,z_2,-y/c,z_3-\bar\tau_R y/c).
\end{equation}
Indeed the scalar factor $c$ in four dimensions contributes $c^4$
and $(\det L_R)^2=c^2$. The density transformation consequently gives
\begin{equation}\label{eq:v37-return-gaussian-agreement}
 c\,g_{\Sigma_R}(z)g_{\sigma_{A,R}^2}(y-\beta_Rz)
   =c\,g_{\Omega_R}(z,y)
   =c^{-2}g_{D_R}\bigl(c^{-1/2}L_R^{-1}(z,y)\bigr).
\end{equation}
For $n=mc+\sqrt m\,y$, the centered return coordinates divided by
$\sqrt m$ are exactly
$(z_1,z_2,-y/c,z_3-\bar\tau_Ry/c)$.
Since $n/m\to c$ on bounded $y$ windows, the argument and prefactor
on the right of \eqref{eq:v37-return-gaussian-agreement} are precisely
those obtained by summing a putative $n^{-2}$ four-coordinate local
law over $\sqrt m$ consecutive return counts. This calculation checks
the normalization; it is not an inference of that pointwise law.

\paragraph{The remaining inversions.}
Theorem~\ref{thm:v37-return-window} fixes both displacement coordinates
and the collision count exactly, keeps a roof interval of unscaled
positive length, and resolves the actual return index on its diffusive
scale. It supplies a local constraint for the true four-coordinate
return measure, not merely for the stationary flow. To resolve one
return index one must control occupation frequencies outside the
shrinking neighborhood used in
\eqref{eq:v37-compact-small-occupation}. To resolve a pointwise roof
density one must also control the raw edge correction and the full
roof-frequency complement. Neither inverse is supplied by convergence
of the window measures. The common pointwise correction and the
complete return complement in the original raw theorem are retained
as distinct analytical requirements.
